% =============================================================
% Capitolo: Retrieval-Augmented Generation (RAG) e Memoria Globale
% Progetto: MCP_ECOM — Agente e-commerce basato su MCP
% =============================================================

\chapter{Architettura RAG e Sistema di Memoria Globale}
\label{chap:rag}

% ─────────────────────────────────────────────────────────────
\section{Introduzione al Modello RAG nell'E-commerce}
\label{sec:rag-intro}
% ─────────────────────────────────────────────────────────────

Nel contesto del progetto MCP\_ECOM, l'architettura \textbf{RAG (Retrieval-Augmented Generation)} svolge una funzione cruciale: superare i limiti delle interrogazioni puramente \textit{live} alle API di eBay, dotando l'agente di una \textbf{Memoria Globale} persistente.

Nelle classiche pipeline e-commerce, ogni ricerca è effimera: terminata l'elaborazione, i dati vengono scartati e l'integrazione di contesto storico è tipicamente nulla. Per un agente decisionale autonomo, invece, è fondamentale poter ricordare e recuperare prodotti visti in passato, specialmente articoli interessanti (ma magari sfuggiti alle query attuali) o profili di venditori sui quali si è già calcolata l'affidabilità. 

Il sistema RAG si fa carico di:
\begin{itemize}
    \item \textbf{Ingestion (Indicizzazione)}: Memorizzare i metadati e i contenuti semantici dei prodotti visualizzati in un Vector Database locale.
    \item \textbf{Retrieval (Recupero)}: Identificare e pescare asincronamente risultati pregressi pertinenti alla nuova richiesta dell'utente tramite vettori densi e sparsi.
    \item \textbf{Merging \& Reranking}: Fondere le risposte dell'API in tempo reale (Live eBay) con i ricordi semantici, applicando controlli di aderenza ai vincoli e bonus di rilevanza.
\end{itemize}

% ─────────────────────────────────────────────────────────────
\section{Architettura Ibrida: Qdrant e FastEmbed}
\label{sec:rag-arch}
% ─────────────────────────────────────────────────────────────

Il nucleo di memorizzazione semantica del RAG è appoggiato a \textbf{Qdrant}, un Vector Search Engine ottimizzato, in esecuzione in modalità file-system locale (\texttt{path="qdrant\_storage"}).

Per massimizzare la precisione del recupero, il sistema adotta un modello di \textbf{Hybrid Search} (Ricerca Ibrida), affiancando il tradizionale embedding denso all'embedding sparso basato su frequenze documentali (BM25 locale).

\subsection{Dense Embedding vs Sparse Embedding}
\label{subsec:rag-embeddings}

Ogni articolo inserito in memoria attraversa una duplice codifica:
\begin{description}
    \item[Vettori Densi (Dense Vectors)] Rappresentano il significato semantico latente. Utilizzano un modello \textit{Sentence Transformer} locale con dimensione vettoriale $d = 384$. La misurazione della similarità avviene tramite distanza coseno. Ottimi per gestire sinonimi o concetti latenti.
    \item[Vettori Sparsi (Sparse Vectors)] Implementati tramite la libreria \texttt{fastembed} con modello \texttt{Qdrant/bm25}. Catturano l'esatta presenza sintattica delle parole e la loro rilevanza tramite ponderazione \textit{Term Frequency - Inverse Document Frequency (TF-IDF)}. Sono eccellenti per intercettare termini specifici, come codici prodotto, esatti nomi di brand o numeri di modello (es. "GTX 1660").
\end{description}

\subsection{Reciprocal Rank Fusion (RRF)}
\label{subsec:rag-rrf}

La fase di recupero (query) sfrutta nativamente la pipeline \texttt{Prefetch} di Qdrant unita al meccanismo \textbf{Reciprocal Rank Fusion (RRF)}. Piuttosto che calcolare medie ponderate tra gli score di natura differente, l'RRF riordina basandosi sui rank di ciascuna coda:
\begin{equation}
    RRFScore(d) = \sum_{q \in Queries} \frac{1}{k + rank(d, q)}
\end{equation}
Qdrant esegue un prefetch sia sui vettori densi (limite empirico configurato a $2k$) sia sui vettori sparsi (limite $2k$), fondendo le due distribuzioni prima di restituire il set top-$k$ prefiltrato.

% ─────────────────────────────────────────────────────────────
\section{Ingestion: Costruzione del Paylod Documentale}
\label{sec:rag-ingestion}
% ─────────────────────────────────────────────────────────────

L'indicizzazione degli articoli (\textit{index\_search\_items}) avviene proattivamente durante lo scraping o l'estrazione dai risultati d'asta. 
Affinché i calcoli semantici producano vettori densamente informativi, è essenziale che il testo di input aggreghi i metadati essenziali dell'articolo (\textit{Title Enrichement}).

\begin{verbatim}
enrich_text = f"{title}"
if item.get("brand"):
    enrich_text += f" brand:{item.get('brand')}"
if item.get("condition"):
    enrich_text += f" condition:{item.get('condition')}"
\end{verbatim}

\subsection{Generazione deterministica dei Point ID}
La gestione degli identificativi in Qdrant accetta limitatamente \textit{str UUID} o Interi. Per garantire l'idempotenza durante gli \textit{upsert} (inserimenti/aggiornamenti evitano duplicati qualora il sistema visiti più volte lo stesso annuncio), il sistema converte l'\texttt{ebay\_id} univoco in UUID attraverso una funzione di hashing MD5 standard:

\begin{verbatim}
import hashlib
h = hashlib.md5(doc_id_str.encode("utf-8")).hexdigest()
point_id = str(uuid.UUID(h))
\end{verbatim}

I metadati (comprensivi di prezzo originale, rating del venditore, condition e url) vengono isolati da artefatti e serializzati tramite parser \textit{Pydantic-NumPy Sanitization} (\texttt{sanitize\_scalars}) prima di essere iniettati nel payload di Qdrant.

% ─────────────────────────────────────────────────────────────
\section{Retrieval: Integrazione nella Search Pipeline}
\label{sec:rag-pipeline}
% ─────────────────────────────────────────────────────────────

All'arrivo di una richiesta utente, la pipeline esegue il cosiddetto processo \textbf{Phantom Fetch}. Onde evitare l'innalzamento dei cicli di clock e latenza in attesa alle chiamate Network verso eBay, l'estrapolazione in Qdrant è confinata in un Thread asincrono parallelizzato al fetching esterno (\texttt{asyncio.create\_task(asyncio.to\_thread(...))}).

\subsection{Query Expansion (Espansione Autonoma)}
Un micro-servizio LLM innesca dinamicamente l'\textbf{Expand Query}: la stringa primigenia dell'utente (es. \textit{"scarpa da tennis azzurra"}) viene de-codificata, privata di stop words e allargata ai sinonimi semantici, potenziando drasticamente il recall del Vector DB (\textit{"scarpa sneaker calzatura tennis corsa azzurra blu ciano"}).

\subsection{Logica di Fusione e Filtraggio Constraint (\_merge\_hybrid\_results)}
Una volta terminate le due query concorrenti (eBay Live + Qdrant Memoria), i due listati subiscono una brutale intersecazione logica.

Se un annuncio "fantasma" arriva dalla Memoria, subisce restrizioni fortissime (Hard Filters) prima di avere accesso all'orologio del Reranker finale:
\begin{itemize}
    \item \textbf{Overlap Soggettivo}: Almeno il 50\% delle parole chiave utili nella stringa della Query originale devono figurare nel titolo.
    \item \textbf{Match Esatto di Prodotto}: Se l'LLM ha dedotto un \texttt{product} cardine (\textit{product constraint extraction}), il lemma base (depauperato da inflessioni plurali) deve obbligatoriamente persistere nell'annuncio memoria.
    \item \textbf{Filtri Prezzo / Condizione Logici}: Analisi condizionale sui vincoli \texttt{<=, >=, between} estratti preventivamente, tagliando ogni ramo inosservante senza pietà.
\end{itemize}

\subsection{Verifica in tempo reale API (Dead Listing check)}
\label{subsec:rag-validation}
Una forte vulnerabilità delle memorie vettoriali su scenari market asincroni (dove un lotto può scadere arbitrariamente) consiste nell'esporre all'utente un ricordo, sebbene logicamente congruente, oramai deceduto materialmente.

Per salvaguardare l'accuratezza, la pipeline setaccia gli articoli provenienti unicamente dal RAG (che \textit{non} sono accoppiati da rinvenimenti Live d'eBay). Su questa nicchia viene imbastito un blocco asincrono massivo (\texttt{asyncio.gather}) puntando singolarmente all'API \texttt{get\_item\_details(ebay\_id)}. Se l'endpoint eBay dà responso fallimentare o nullo, l'identificativo viene tacciato come \textit{Dead ID} e l'oggetto estirpato dalla fusione semantica. Ove vitale si beneficia per aggiornare metadati instabili come Pricing o Spedizioni, fornendo freschezza istantanea d'informazione (\textit{Data Enrichment}).

\subsection{Calcolazione dell'Hybrid Bonus}
Al termine degli smistamenti i due universi (Live e Memoria) si combinano. A qualunque listing che figura in \textit{entrambe i listati} (essendo scovato sul Web adesso ma risultante egualmente da profitti semantici incassati sui vettori storici passati) viene insignito un \textbf{hybrid bonus} pari a 0.15. 
Questo segnale rinforzato decreta un indizio empirico di importanza colossale e si appoggia alla logica base del Sort Array tramite inverse rank ponderale:
\begin{equation}
    SortWeight = \frac{1}{\rho_{\text{eBay}} + 10} + \frac{0.5}{\rho_{\text{Memoria}} + 10} + Bonus_{\text{Hybrid}}
\end{equation}

% ─────────────────────────────────────────────────────────────
\section{Sfide Architetturali, Limitazioni Server e Future Iterazioni}
\label{sec:rag-critics}
% ─────────────────────────────────────────────────────────────

Sebbene l'architettura si sia dimostrata una cerniera innovativa tra la ritenzione d'informazione ed un workflow ReAct flessibile, solleva delle problematiche architetturali in ambito esecutivo:

\begin{description}
    \item[Concorrenza e Locks di Qdrant Local] Operando Qdrant su filesystem in ambiente \textit{in-process} originario tramite l'engine SQLite/RocksDB sottostante si manifestano crash catastrofici non appena \texttt{Uvicorn} effettui uno scale verticale (\textit{Workers $>$ 1}). Avendo i Thread Worker una spartizione parallela, incappano in un Lock file fatale alla cartella "qdrant\_storage". Questo esige l'evaporazione verso un \textit{Microservizio Docker Qdrant standalone} in caso di migrazione in prod.
    
    \item[Rate Limiting d'Arricchimento API] Il sistema \texttt{gather} adibito alla convalida in Real-Time (\textit{Deep-Check eBay Details}) per i filettaggi della memoria avvia dozzine di chiamate HTTP sincrone concorrenti. Questa voracità network disarmante, senza adeguati semafori asincroni (es. \texttt{asyncio.Semaphore(5)}), incappa in Rate-Limit severi e \textit{IP Blocking / HTTP 429} dalle barricate d'eBay qualora il traffico salisse drasticamente in un'unica instanza.
    
    \item[Memoria Iniziale ad Aggancio Lazy] L'inizializzazione del modello Transformer di \texttt{fastembed} in \texttt{qdrant\_store.py} viene demandata brutalmente \textit{Lazy} alla prima istanza di richiesta. Tale design posticipa uno spike spropositato I/O di memoria Ram alla prima query che potrebbe stallare i pool asincroni.
\end{description}
    
Questa architettura RAG rappresenta una svolta vitale nell'ibridazione decisionale: l'unione rigorosa di estrazione Live transitoria con un motore latente e stratificato ad Hybrid Search vettoriale, capace di rielaborare, autogiustificarsi e ricalibrare i falli.
